\chapter*{Preface}
\addcontentsline{toc}{chapter}{Preface}
\markboth{Preface}{Preface}

Amateur (``ham'') radio sits at a rare intersection: it is simultaneously a
hands-on engineering discipline, an experimental science, a public-service
tradition, and a global community held together by little more than
electromagnetic waves and a shared curiosity about how things work. Every
licensed amateur, anywhere in the world, has had to learn the same
underlying physics --- Ohm's law does not change at a border, and a
half-wave dipole behaves identically whether it is strung between two trees
in Lisbon, Nairobi, or Auckland.

This book was written to teach that underlying physics and engineering as
completely and rigorously as a motivated newcomer needs, independent of
where they happen to live. National licensing authorities differ enormously
in their exam formats, their band allocations, their power limits, and the
paperwork required to get on the air --- and this book intentionally does
\emph{not} try to teach any single country's regulations. Instead, it
teaches the physics, mathematics, electronics, antenna theory, propagation,
and operating practice that sit underneath \emph{every} national exam
syllabus, so that the material remains useful and accurate regardless of
which administration issues your callsign. Where a topic is genuinely
international --- the ITU Region system, the Q-code, the NATO phonetic
alphabet, Morse operating abbreviations --- it is covered in full, because
these are shared worldwide conventions rather than local law.

The book is organized in five parts of increasing sophistication, loosely
mirroring the tiered structure used by most national licensing systems
(an entry-level tier, an intermediate tier, and a full-privileges tier),
so that a reader preparing for a foundation-level exam can stop after
Part~II with a solid, self-sufficient understanding, while a reader aiming
for the highest available privileges can continue through the treatment of
active devices, transmission line theory, antenna engineering, and
propagation in Parts~III and~IV. Part~V turns from theory to practice:
operating procedure, radio etiquette, digital modes, and emergency
communications, all described in terms general enough to apply to any
amateur service worldwide.

Because exam questions ultimately test whether you can \emph{use} a
formula, not just recite it, every chapter includes worked examples,
exam-style tips, and end-of-chapter review questions. A comprehensive
formula appendix (Appendix~A) collects every equation in the book in one
place for quick reference and last-minute revision, alongside a table of
physical constants, SI prefixes, and unit conversions (Appendix~B) and a
glossary of terms (Appendix~C).

Before you use this book to sit a real examination, confirm the specific
band plan, power limits, callsign format, and administrative procedure
that apply in your own country with your national telecommunications
regulator or IARU member society --- those details are the one part of
becoming a radio amateur that genuinely does depend on where you live.

Good luck, and 73.
